% Gerado por generate_scientific_figures.py: cada fórmula aparece abaixo do algoritmo.
% Matemática em \footnotesize, sem ampliação; somente o somatório de Takano usa escala 0.95.
\begingroup
\small
\setlength{\tabcolsep}{4pt}
\setlength{\LTcapwidth}{\linewidth}
\setlength{\LTpre}{\smallskipamount}
\setlength{\LTpost}{\smallskipamount}
\renewcommand{\arraystretch}{1.10}
\begin{longtable}{@{}>{\raggedleft\arraybackslash}p{1.8cm} p{\dimexpr\linewidth-1.8cm-2\tabcolsep\relax}@{}}
\caption[Catálogo de fórmulas analíticas — Gemini. As expansões de arco-tangente são exibidas explicitamente em forma de somatório.]{\footnotesize\sffamily\bfseries Catálogo de fórmulas analíticas — Gemini. As expansões de arco-tangente são exibidas explicitamente em forma de somatório.}\label{tab:formula-catalog-gemini}\\
\toprule
Prioridade & Método e fórmula (a expressão fica abaixo do nome)\\
\midrule
\endfirsthead
\toprule
Prioridade & Método e fórmula (a expressão fica abaixo do nome)\\
\midrule
\endhead
\midrule
\multicolumn{2}{r}{\footnotesize Continua na próxima página.}\\
\endfoot
\bottomrule
\endlastfoot
1 & \texttt{machin}\\*[2pt]
\multicolumn{2}{@{}>{\centering\arraybackslash}p{\linewidth}@{}}{\mbox{\footnotesize $\displaystyle \pi=16\arctan(1/5)-4\arctan(1/239)$}}\\*[3pt]
\multicolumn{2}{@{}>{\centering\arraybackslash}p{\linewidth}@{}}{\mbox{\footnotesize $\displaystyle 16\sum_{k=0}^{\infty}\frac{(-1)^k}{(2k+1)5^{2k+1}}-4\sum_{k=0}^{\infty}\frac{(-1)^k}{(2k+1)239^{2k+1}}$}}\\[5pt]
2 & \texttt{euler}\\*[2pt]
\multicolumn{2}{@{}>{\centering\arraybackslash}p{\linewidth}@{}}{\mbox{\footnotesize $\displaystyle \pi=2\sum_{k=0}^{\infty}\frac{k!}{(2k+1)!!}$}}\\[5pt]
3 & \texttt{newton}\\*[2pt]
\multicolumn{2}{@{}>{\centering\arraybackslash}p{\linewidth}@{}}{\mbox{\footnotesize $\displaystyle \pi=3\sum_{k=0}^{\infty}\frac{(2k)!}{16^k(k!)^2(2k+1)}$}}\\[5pt]
4 & \texttt{nilakantha}\\*[2pt]
\multicolumn{2}{@{}>{\centering\arraybackslash}p{\linewidth}@{}}{\mbox{\footnotesize $\displaystyle \pi=3+4\sum_{k=1}^{\infty}\frac{(-1)^{k+1}}{2k(2k+1)(2k+2)}$}}\\[5pt]
5 & \texttt{ramanujan}\\*[2pt]
\multicolumn{2}{@{}>{\centering\arraybackslash}p{\linewidth}@{}}{\mbox{\footnotesize $\displaystyle \pi=\frac{9801}{2\sqrt{2}\sum_{k=0}^{\infty}\frac{(4k)!(1103+26390k)}{(k!)^4 396^{4k}}}$}}\\[5pt]
\end{longtable}
\endgroup
